% Part: incompleteness
% Chapter: representability-in-q
% Section: beta-function

\documentclass[../../../include/open-logic-section]{subfiles}

\begin{document}

% SEGMENT OLP-0292-S01

\olfileid{inc}{req}{bet}
\olsection{பீட்டா சார்புத் துணைத்தேற்றம்}


கூட்டல், பெருக்கல், $\Char{=}$ ஆகியவை கிடைக்கும்போது முதன்மை
மீள்வரையறையைச் செயல்படுத்த முடியும் என்று காட்டுவதற்கு, தொடர்களைக் கையாளும்
சார்புகளை உருவாக்க வேண்டும். (அடுக்கேற்றமும் நம்மிடம் இருந்தால், நமது பணி
எளிதாக இருக்கும்.) முதன்மை மீள்வரையறை நம்மிடம் இருந்தபோது, ``$n$-ஆம்
பகா எண்'' போன்றவற்றை வரையறுத்து, ஓரளவு நேரடியான குறியீடாக்கத்தைத்
தேர்ந்தெடுக்க முடிந்தது. ஆனால் இங்கு முதன்மை மீள்வரையறை நம்மிடம்
இல்லை---உண்மையில் சிறுமமாக்கலைப் பயன்படுத்தி முதன்மை மீள்வரையறையைச் செய்ய
முடியும் என்று காட்ட விரும்புகிறோம்---ஆகவே இன்னும் நுணுக்கமாகச் செயல்பட
வேண்டும்.

\begin{lem}
\ollabel{lem:beta}
ஒவ்வொரு தொடர் $a_0$, \dots,~$a_n$ க்கும், ஒவ்வொரு $i \le n$ க்கும்
$\beta(d,i) = a_i$ என அமையுமாறு ஓர் எண்~$d$ இருக்கக்கூடிய
$\beta(d,i)$ என்ற சார்பு ஒன்று உள்ளது. மேலும், அடிப்படைச் சார்புகளிலிருந்து
சேர்ப்பு மற்றும் ஒழுங்கான சிறுமமாக்கல் ஆகியவற்றை மட்டும் பயன்படுத்தி
$\beta$ ஐ வரையறுக்க முடியும்.
\end{lem}

$d$ ஆனது தொடர் $\tuple{a_0, \dots, a_n}$ ஐக் குறியீடாக்குவதாகவும்,
$\beta(d,i)$ அதன் $i$-ஆம் உறுப்பைத் திருப்பித்தருவதாகவும் கருதுக.
(இந்தக் ``குறியீடாக்கம்'' நாம் ஏற்கெனவே அறிந்த பகா எண் அடுக்குக்
குறியீடாக்கத்தைப் \emph{பயன்படுத்துவதில்லை} என்பதைக் கவனிக்க!) துணைத்தேற்றம்
மிகக் குறைந்த அளவிலானதே; தொடர்களைத் தொடரிணைக்கவோ, உறுப்புகளை இறுதியில்
சேர்க்கவோ முடியும் என்றுகூட அது கூறவில்லை; சேர்ப்பு மற்றும் ஒழுங்கான
சிறுமமாக்கலால் வரையறுக்கத்தக்க சார்புகளைப் பயன்படுத்தி $a_0$, \dots,~$a_n$
ஆகியவற்றிலிருந்து $d$ ஐக் \emph{கணிக்க} முடியும் என்றுகூட அது கூறவில்லை.
ஒவ்வொரு தொடரும் ``குறியீடாக்கப்படும்'' வகையில் ஒரு ``குறிவிலக்கச்'' சார்பு
உள்ளது என்று மட்டுமே அது கூறுகிறது.

$\beta$ என்ற குறியீட்டைப் பயன்படுத்துவது கோடலிடமிருந்து வந்தது. மீண்டும்
கூறினால், பொருத்தமான~$\beta$ ஒன்றை வெளித்தோற்றத்தில் வரம்புடைய வளங்களை
மட்டும் கொண்டு, அதாவது சேர்ப்பு மற்றும் சிறுமமாக்கலை மட்டும் கொண்டு,
வரையறுப்பதே துணைத்தேற்றத்தை நிறுவுவதிலுள்ள கடினமான பகுதி---ஆனால் கூட்டல்,
பெருக்கல், $\Char{=}$ ஆகியவற்றைப் பயன்படுத்த நமக்கு அனுமதி உள்ளது.
இத்துணைத்தேற்றத்தை நிறுவப் பல வழிகள் உள்ளன; அவற்றில் மிகத் தெளிவான
ஒன்று, சுன்சுவின் தேற்றம் (மரபாக, ``சீன மீதித் தேற்றம்'') எனப்படும்
எண்-கோட்பாட்டு உண்மையைப் பயன்படுத்திய கோடலின் மூல முறையே ஆகும்.

\begin{defn}
இயல் எண்கள் $a$,~$b$ ஆகியவற்றின் மீப்பெருப் பொதுவகுத்தி~$1$ எனில்,
எனில் மட்டுமே, அவை \emph{சார்பகா} எண்கள்; வேறுவிதமாகக் கூறினால்,
வேறு எந்த வகுத்தியையும் அவை பொதுவாகக் கொண்டிருக்கவில்லை.
\end{defn}

\begin{defn}
இயல் எண்கள் $a$,~$b$ ஆகியவை \emph{மட்டு~$c$ க்கு ஒருங்கிசைவு உடையவை},
$a \equiv b \mod c$, எனில், எனில் மட்டுமே, $c \mid (a-b)$; அதாவது,
$c$ ஆல் வகுக்கப்படும்போது $a$,~$b$ இரண்டும் ஒரே மீதியைக் கொண்டுள்ளன.
\end{defn}

சுன்சுவின் தேற்றம் இதுதான்:
\begin{thm}
$x_0$, \dots,~$x_n$ ஆகியவை (சோடிசோடியாகச்) சார்பகா எண்கள் என்று கொள்க.
$y_0$, \dots,~$y_n$ ஏதேனும் எண்களாக இருக்கட்டும். அப்போது பின்வருமாறு
அமையும் ஓர் எண் $z$ உள்ளது:
\begin{align*}
z & \equiv y_0 \mod x_0 \\
z & \equiv y_1 \mod x_1 \\
& \vdots  \\
z & \equiv y_n \mod x_n.
\end{align*}
\end{thm}

சுன்சுவின் தேற்றத்தை நாம் பின்வருமாறு பயன்படுத்துவோம்: $x_0$,
\dots,~$x_n$ ஆகியவை முறையே $y_0$, \dots,~$y_n$ ஆகியவற்றைவிடப்
பெரியவை எனில், தொடர் $\tuple{y_0, \dots,y_n}$ ஐக் குறியீடாக்க $z$ ஐ
எடுத்துக்கொள்ளலாம். $y_i$ ஐ மீட்க, $z$ ஐ~$x_i$ ஆல் வகுத்து மீதியை
மட்டும் எடுத்தால் போதும். இக்குறியீடாக்கத்தைப் பயன்படுத்துவதற்கு,
$x_0$, \dots,~$x_n$ க்குப் பொருத்தமான மதிப்புகளைக் கண்டறிய வேண்டும்.

இவ்வகையில் சில கவனிப்புகள் நமக்கு உதவும். $y_0$, \dots,~$y_n$
தரப்பட்டால், பின்வருமாறு கொள்க:
\begin{align*}
j &= \max(n, y_0 + 1, \dots, y_n + 1), \\
m &= \lcm(1,\dots,j),
\end{align*}
மேலும் பின்வருமாறு கொள்க:
\begin{align*}
x_0 & = 1 + m \\
x_1 & = 1 + 2 \cdot m \\
x_2 & = 1 + 3 \cdot m \\
& \vdots  \\
x_n & = 1 + (n+1) \cdot m
\end{align*}
அப்போது இரண்டு உண்மைகள் நிலவுகின்றன:
\begin{enumerate}
\item\ollabel{rel-prime} $x_0,\dots,x_n$ ஆகியவை சார்பகா எண்கள்.
\item\ollabel{less} ஒவ்வொரு $i$ க்கும், $y_i < x_i$.
\end{enumerate}
\olref{rel-prime} மெய் என்று காண, $p$ ஒரு பகா எண் என்றும்,
$p \mid x_i$, $p \mid x_k$ என்றும் இருந்தால், $p \mid 1 + (i+1) m$
மற்றும் $p \mid 1 + (k+1) m$ என்பதைக் கவனிக்க. ஆனால் அப்போது அவற்றின்
வேறுபாட்டை $p$ வகுக்கிறது:
\[
(1 + (i+1)m) - (1+ (k+1)m) = (i-k) m.
\]
$p$ ஆனது $1 + (i+1)m$ ஐ வகுப்பதால், அது $m$ ஐயும் வகுக்க முடியாது
(அவ்வாறு வகுத்தால், முதல் வகுத்தல்~$1$ ஐ மீதியாக விட்டிருக்கும்). எனவே
$p$ ஆனது $i-k$ ஐ வகுக்கிறது; ஏனெனில் $p$ ஆனது $(i-k)m$ ஐ வகுக்கிறது. ஆனால்
$\left|i-k\right|$ அதிகபட்சம்~$n$; மேலும் $j \geq n$ எனத்
தேர்ந்தெடுத்துள்ளோம். ஆகவே மீண்டும் ஒரு முரணாக $p \mid m$ என்பது இதிலிருந்து
பின்வருகிறது. எனவே $x_i$,~$x_k$ இரண்டையும் வகுக்கும் பகா எண் எதுவும்
இல்லை. கூறு~\olref{less} எளிதானது: $y_i < j \leq m < x_i$.

இப்போது $\beta$ சார்புத் துணைத்தேற்றத்தை நிறுவுவோம். $0$, அடுத்தெண்,
கூட்டல், பெருக்கல், $\Char{=}$, வீழ்த்தல்கள், மேலும் இவற்றிலிருந்து
சேர்ப்பு மற்றும் ஒழுங்கான சார்புகளுக்குப் பயன்படுத்தப்படும் சிறுமமாக்கல்
ஆகியவற்றால் வரையறுக்கப்படும் எந்தச் சார்பையும் பயன்படுத்தலாம் என்பதை
நினைவுகூர்க. ஒரு தொடர்பின் சிறப்பியல்புச் சார்பு இவ்வாறு வரையறுக்கத்தக்கது
எனில், அத்தொடர்பையும் பயன்படுத்தலாம். முன்புபோலவே இத்தொடர்புகள் பூலியச்
சேர்க்கைகள் மற்றும் வரம்பிட்ட அளவையிடலின்கீழ் மூடியவை என்று காட்டலாம்;
எடுத்துக்காட்டாக:
\begin{align*}
\fn{not}(x) & \defis \Char{=}(x,0)\\
\bmin{x \leq z}{R(x,y)} & \defis \umin{x}{(R(x,y) \lor x = z)}\\
\bexists{x \leq z}{R(x,y)} & \defiff R(\bmin{x \leq z}{R(x,y)}, y)
\end{align*}
பின்னர் பின்வருவன அனைத்தும் முதன்மை மீள்வரையறை இல்லாமலேயே
வரையறுக்கத்தக்கவை என்றும் காட்டலாம்:
\begin{enumerate}
\item சோடியாக்கச் சார்பு, $J(x,y) = \frac{1}{2}[(x+y)(x+y+1)] + x$;
% இங்கு என்ன நடக்கிறது என்பதை மேலும் விளக்கலாம்; சற்று குழப்பமாக உள்ளது.
\item வீழ்த்தல் சார்புகள்
\begin{align*}
K(z) & = \bmin{x \leq z}{\bexists{y \leq z}{z = J(x,y)}},\\
L(z) & = \bmin{y \leq z}{\bexists{x \leq z}{z = J(x,y)}};
\end{align*}
\item குறைவு தொடர்பு $x < y$;
\item மீதியின்றி வகுபடும் தொடர்பு $x \mid y$;
% \item $x \tsub y$
% \item $\fn{Prime}(x)$
% \item $p$ பகா எண் எனக் கொண்டு, ``$x$ ஆனது $p$ இன் ஓர் அடுக்கு'' என்ற தொடர்பு:
% \[
% \bforall{y \leq x}{(y \mid x \lif y = 1 \lor y = x)}.
% \]
\item $y$ ஐ~$x$ ஆல் வகுக்கும்போது கிடைக்கும் மீதியைத் திருப்பித்தரும்
  சார்பு~$\fn{rem}(x,y)$.
\end{enumerate}
இப்போது பின்வருமாறு வரையறுக்க:
\begin{align*}
\beta^*(d_0,d_1,i) & = \fn{rem}(1+(i+1) d_1,d_0) \text{ மற்றும்}\\
\beta(d,i) & = \beta^*(K(d),L(d),i).
\end{align*}
இதுவே நமக்குத் தேவையான சார்பு. மேலுள்ளவாறு $a_0,\dots,a_n$
தரப்பட்டால்,
\[
j = \max(n,a_0+1,\dots,a_n+1),
\]
என்க; மேலும் $d_1 = \lcm(1,\dots,j)$ என்க. மேலுள்ள
\olref{rel-prime} இன்படி, $1+d_1$, $1+2 d_1$, \dots,
$1+(n+1) d_1$ ஆகியவை சார்பகா எண்கள் என்றும், \olref{less} இன்படி
அவை அனைத்தும் முறையே $a_0,\dots,a_n$ ஆகியவற்றைவிடப் பெரியவை என்றும்
அறிகிறோம். சுன்சுவின் தேற்றத்தின்படி, ஒவ்வொரு~$i$ க்கும் பின்வருமாறு
அமையும் மதிப்பு~$d_0$ ஒன்று உள்ளது:
\[
d_0 \equiv a_i \mod (1+(i+1)d_1)
\]
ஆகவே ($d_1$ ஆனது~$a_i$ ஐவிடப் பெரியது என்பதால்),
\[
a_i = \fn{rem}(1+(i+1)d_1,d_0).
\]
$d = J(d_0,d_1)$ என்க. அப்போது ஒவ்வொரு $i \le n$ க்கும்,
\begin{align*}
\beta(d,i) & =  \beta^*(d_0,d_1,i) \\
& =  \fn{rem}(1+(i+1) d_1,d_0) \\
& =  a_i
\end{align*}
என அமைகிறது; இதுவே நமக்குத் தேவை. இத்துடன் $\beta$ சார்புத்
துணைத்தேற்றத்தின் நிறுவல் முடிகிறது.

\begin{prob}
  தொடர்புகள் $x < y$, $x \mid y$, சார்பு~$\fn{rem}(x,y)$ ஆகியவற்றை
  முதன்மை மீள்வரையறை இல்லாமல் வரையறுக்க முடியும் என்று காட்டுக. $0$,
  அடுத்தெண், கூட்டல், பெருக்கல், $\Char{=}$, வீழ்த்தல்கள், வரம்பிட்ட
  சிறுமமாக்கல் மற்றும் அளவையிடல் ஆகியவற்றைப் பயன்படுத்தலாம்.
\end{prob}

\end{document}
